% ============================================================
%  SEÇÃO 2 — REFERENCIAL TEÓRICO
%
%  O referencial teórico mostra que você conhece a área.
%  Organize por TEMAS, não por autores.
%  Estrutura sugerida:
%    2.1 Conceito principal do seu tema
%    2.2 Tecnologia / abordagem central
%    2.3 Trabalhos relacionados (o que já foi feito)
%
%  Dica: mínimo de 15 referências citadas no texto (Ficha I).
%  Tamanho sugerido: 8 a 15 páginas
% ============================================================

\chapter{REFERENCIAL TEÓRICO}
\label{cap:referencial}

% ----------------------------------------------------------
%  COMO CITAR com abntex2 + biblatex-abnt
%
%  \cite{chave}         $\rightarrow$ (AUTOR, ano)       — citação indireta
%  \textcite{chave}     $\rightarrow$ Autor (ano)        — citação indireta no texto
%  \cite[p.~10]{chave}  $\rightarrow$ (AUTOR, ano, p. 10) — com número de página
%  \apud{chave}{chave2} $\rightarrow$ (A apud B, ano)    — citação de citação (evitar)
%
%  A chave é gerada automaticamente pelo Zotero + Better BibTeX.
%  Exemplo de chave: silva2022sistemasweb
% ----------------------------------------------------------


% ----------------------------------------------------------
% 2.1 PRIMEIRO TEMA
% Nome este subseção com o conceito central do seu trabalho
% Exemplos: "Sistemas de Gestão", "Inteligência Artificial",
%           "Arquitetura de Microsserviços", "LGPD"
% ----------------------------------------------------------

\section{[Nome do Primeiro Tema]}

% %->  ESCREVA AQUI a fundamentação do primeiro tema.
%     Defina o conceito, apresente as principais visões
%     teóricas e cite os autores relevantes.
%
% Exemplo de parágrafo com citação indireta:
% Sistemas de informação podem ser definidos como conjuntos
% organizados de pessoas, hardware, software e redes de
% comunicação que coletam, transformam e distribuem
% informações em uma organização \cite{laudon2021}.
%
% Exemplo de citação no texto:
% Segundo \textcite{pressman2021}, a engenharia de software...

[Escreva aqui o desenvolvimento do primeiro tema. Cite ao menos 3 referências.]

\subsection{[Subtema 1.1 — se necessário]}

% %->  Use subseções para aprofundar aspectos específicos.
%     Evite subseções com menos de 3 parágrafos.

[Desenvolvimento do subtema, se necessário.]


% ----------------------------------------------------------
% 2.2 SEGUNDO TEMA
% Foque na tecnologia, método ou abordagem que você vai usar
% ----------------------------------------------------------

\section{[Nome do Segundo Tema]}

% %->  ESCREVA AQUI o segundo bloco temático.
%     Este geralmente aborda a tecnologia ou metodologia
%     central do seu projeto.
%
% Exemplos de temas possíveis:
%  - Desenvolvimento de aplicações web (React, Node.js, etc.)
%  - Machine Learning e Redes Neurais
%  - Arquitetura de Microsserviços
%  - Segurança em APIs REST
%  - Design Science Research como metodologia

[Escreva aqui o segundo tema. Cite ao menos 3 referências diferentes das anteriores.]

\subsection{[Subtema 2.1 — se necessário]}

[Desenvolvimento do subtema, se necessário.]


% ----------------------------------------------------------
% 2.3 TERCEIRO TEMA (opcional — adicione se necessário)
% ----------------------------------------------------------

\section{[Nome do Terceiro Tema]}

% %->  Adicione quantas seções forem necessárias para cobrir
%     todos os conceitos relevantes do seu projeto.

[Escreva aqui o terceiro tema.]


% ----------------------------------------------------------
% 2.4 TRABALHOS RELACIONADOS
% O que já foi feito na área? O que é diferente no seu trabalho?
% ----------------------------------------------------------

\section{Trabalhos Relacionados}

% %->  Apresente 3 a 5 trabalhos (artigos, TCCs, dissertações)
%     com proposta semelhante à sua.
%     Para cada um, descreva:
%       (a) O que o trabalho fez
%       (b) Quais resultados obteve
%       (c) O que você vai fazer diferente
%
% Dica: busque no Google Scholar, IEEE Xplore e BDTD.
% Use a busca: "seu tema" AND "sistemas de informação"

Esta seção apresenta trabalhos com objetivos similares aos deste projeto, evidenciando as contribuições originais propostos.

% Trabalho 1
\textcite{[chave1]} [descreva o que o trabalho fez e o que encontrou].

% Trabalho 2
\textcite{[chave2]} [descreva o que o trabalho fez e o que encontrou].

% Trabalho 3
\textcite{[chave3]} [descreva o que o trabalho fez e o que encontrou].

% Diferencial do seu trabalho
Diferentemente dos trabalhos citados, este projeto [explique o que torna seu trabalho diferente ou complementar aos anteriores].


% ----------------------------------------------------------
%  EXEMPLO COMPLETO DE PARÁGRAFO BEM REFERENCIADO
%  (você pode apagar este bloco após entender a estrutura)
% ----------------------------------------------------------

%\begin{exemplo}
% Exemplo de parágrafo completo com citações:
%
% A gestão de filas em ambientes de saúde representa um desafio
% recorrente no sistema público brasileiro \cite{ministerio2021}.
% Segundo \textcite{silva2022}, a adoção de sistemas digitais
% pode reduzir o tempo médio de espera em até 40\% em unidades
% básicas de saúde. Complementando esta visão,
% \textcite{costa2023} demonstrou que a usabilidade do sistema
% é fator determinante para a adoção por parte dos usuários.
% Diante disso, este trabalho propõe [sua contribuição].
%\end{exemplo}
